\chapter{The compiler}\label{ch:compiling}

A program reaches native G17 bytes through one of two input routes and a shared backend (\cref{fig:compile-routes}).
\Cref{ch:running-example-through} follows the running example through each stage; the general rules for selection,
allocation and tensor lowering come after it.

\section{Input routes}\label{ch:source-native-program}

Metal source goes first through Apple's front end, which produces AIR, and the project translates AIR into its own
intermediate representation (IR). Python builders construct that IR directly. Both routes then use the same native
backend.

% Figure 15.1: the compiler's two input routes.
\begin{figure}[htbp]
\centering
\begin{tikzpicture}[node distance=5mm and 5mm]
  \node[apple] (msl) {Metal source\\[-1pt]{\footnotesize written outside the project}};
  \node[apple, right=of msl] (fe) {Apple's front end\\[-1pt]{\footnotesize\code{xcrun metal}}};
  \node[apple, right=of fe] (air) {AIR};
  \node[ours, right=of air] (tr) {AIR translation};
  \node[ours, below=10mm of tr] (ir) {The project's IR};
  \node[ours, left=12mm of ir] (py) {Python kernel builders};
  \node[ours, below=of ir, text width=60mm] (be) {Native backend\\[-1pt]{\footnotesize instruction selection, register
    allocation, encoding, tensor lowering}};
  \node[ours, below=of be, text width=60mm] (prog) {Native program\\[-1pt]{\footnotesize G17 bytes and interface description}};
  \node[ours, below left=9mm and -22mm of prog, text width=46mm] (img) {Image author\\[-1pt]{\footnotesize object, metadata, library, archive}};
  \node[nat, below right=9mm and -22mm of prog, text width=46mm] (rt) {Native runtime\\[-1pt]{\footnotesize builds its own resource and launch state}};
  \draw[flow] (msl) -- (fe); \draw[flow] (fe) -- (air); \draw[flow] (air) -- (tr);
  \draw[flow] (tr) -- (ir); \draw[flow] (py) -- (ir);
  \draw[flow] (ir) -- (be); \draw[flow] (be) -- (prog);
  \draw[flow] (prog.south) -- ++(0,-3mm) -| node[pos=0.75, left, tag] {through Metal} (img.north);
  \draw[flow] (prog.south) -- ++(0,-3mm) -| node[pos=0.75, right, tag] {below Metal} (rt.north);
\end{tikzpicture}
\caption{Input routes to a native program. Metal source passes through Apple's front end to AIR, which the project
translates into its own IR; the model kernels are built in that IR directly. Both routes share the native backend,
whose output feeds both execution paths. Grey: Apple's components; blue and red: the project's.}
\label{fig:compile-routes}
\end{figure}


The Metal-source route tests the compiler on programs written outside the project's model builders. In the
reconciled snapshot, 128 sources of the frozen 197-source census compile (MM~\mm{25.196}). On this route only the
translation from source to AIR is Apple's.

The direct-IR route builds the model kernels: the projections, attention, normalization and generation steps. It
does not invoke Apple's Metal front end.

The backend outputs native instruction bytes together with a description of the
program's buffers, bindings and execution requirements. The two execution paths then use it differently:

\begin{itemize}
\item For the Metal path, the image author packages that output into the object, metadata, library and archive that Metal
  loads.
\item The below-Metal runtime takes the native bytes and builds its own measured resource and launch state (\cref{ch:below-metal-native}).
\end{itemize}

Both routes currently use Apple's decoder to check the emitted instructions. That dependency is part of building and
validating a program; the decoder is not involved in delivering the finished program to the GPU.

\section{The running example in the compiler}\label{ch:running-example-through}

The program is the one of \cref{sec:one-program-ir}, compiled and decoded on the CPU. Its bytes are the ones that ran
below Metal, but nothing in this section executes them.

\subsection*{IR}

The IR is SSA over typed values. A function declares its buffers, each with a binding index and an
element type, and tensor operations are first-class.

\begin{lstlisting}[style=py]
import struct
from agxforge.g17 import cc, ir, scanlink

a = ir.Buffer("A", 1, elem=ir.F16)
b = ir.Buffer("B", 2, elem=ir.F16)
c = ir.Buffer("C", 3, elem=ir.F32)
fn = ir.Function("tensor_example", [a, b, c])
bd = ir.Builder(fn, fn.block("entry"))
bd.tensor_matmul(a, b, c, M=17, N=19, K=16)
i = bd.builtin("threadgroup_position_in_grid")
one = struct.unpack("<I", struct.pack("<f", 1.0))[0]
bd.store_at(c, i, bd.fadd(bd.load(c, i, type=ir.F32), bd.const(one)))
bd.ret()
ir.verify(fn)

program = cc.compile_function(fn)    # 1,232 bytes, 101 instructions; SHA-256 8bb67a3c...c667
image = scanlink.author(program)     # object, metadata, library and archive for the Metal path
\end{lstlisting}

\subsection*{Selection}

\code{cc.select} maps each IR operation to a recovered G17 form (\cref{tab:running-example-selection}).

\begin{xltabular}{\linewidth}{@{}>{\hsize=0.71\hsize}X>{\hsize=1.29\hsize}X@{}}
\caption{Selection in the running example: each IR operation and the G17 form it maps to.}\label{tab:running-example-selection}\\
\toprule
IR operation & G17 form \\
\midrule
\endfirsthead
\tablecontinued{2}
\toprule
IR operation & G17 form \\
\midrule
\endhead
\code{tensor\_matmul} & a tensor body generated by \code{agxforge/g17/tlower.py}
(\cref{ch:tensor-lowering}) \\*
\code{builtin("threadgroup\_position\_in\_grid")} & \op{14059} of the threadgroup position in x \\*
\code{load} & \op{12682} \\*
\code{const} & \op{11842} \\*
\code{fadd} & \op{998} \\*
\code{store\_at} & \op{17229} \\*
\code{ret} & \op{684} \\*
\bottomrule
\end{xltabular}

Each of these forms has mapped fields and measured behaviour; selection refuses any other form by name
(\cref{ch:selection-allocation-general}).

\subsection*{Tensor body}

\code{tlower} turns the 17 × 19 × 16 shape into two tile rows and two tile columns, with the second of
each partial:

\begin{itemize}
\item It computes each lane's addresses from the lane index, read with \op{14060} into R0L,
  as row \code{m = 4*(l>>4) + ((l>>1)\&3)} and column \code{col = (l\&8) + 4*(l\&1)}, then \code{m * ld + col} for each
  leading dimension: the B-row packing that \cref{sec:coordinate-systems} explains. The results sit in R1–R8.
\item It builds the partial-tile masks in 16-bit halves (R11H, R12H, R13H), with \op{612} and
  \op{437}, because the masked forms take a 16-bit mask operand.
\item It loads the A fragments into R20–R23 (tile row 0) and R56–R59 (tile row 1, masked), and the B fragments into
  R68–R71 and R72–R75 (the masked second column). Every load publishes scoreboard slot 0.
\item It issues one \op{5107} per output tile, since K = 16 is a single slice.
  Each MMA waits on slot 0.
\item It stores each tile from its accumulator group (R24, R32, R40 and R48, each eight-aligned) with \op{17257}, or \op{17258} where the tile is partial.
\end{itemize}

\subsection*{Allocation}

The allocation follows the constraints of \cref{ch:register-file}:

\begin{itemize}
\item Each accumulator group starts on a multiple of eight, and each 16-bit fragment on a multiple of four
  (\cref{sec:register-classes}).
\item Each fragment is read by two MMAs, kept (0) by the first and released (16) by the second
  (\cref{sec:release-does}). The epilogue's fp32 add releases both its sources, and the store releases its value and index: all
  are last uses.
\item The tensor body names R0–R8, R11H–R13H and R20–R75. The scalar allocator must keep the union of
  every register a tensor body names out of its own choices. The epilogue's registers, R9 and R16–R18, all fall
  outside that set (\code{\_prepare\_tensor\_composition} in \code{agxforge/g17/cc.py}).
\item The register count, 76, is the highest register used plus one. It is reported, not chosen (\cref{sec:capacity-budgets-spills}).
\end{itemize}

\subsection*{Waits}

Each late value is tracked through a scoreboard slot:

\begin{itemize}
\item The threadgroup-position read publishes slot 0.
\item The plain load waits on slot 0 for its index and publishes slot 7.
\item The fp32 add waits on slot 7.
\item No instruction waits on a slot nothing fills, and nothing reads a late value without its wait (\cref{ch:synchronization}).
\end{itemize}

The compiler's checks verify both rules on the emitted bytes, together with the no-read-after-release rule
(\cref{ch:image-result-validated}).

\subsection*{Executor interface}

\code{program.abi()} records what an executor needs:

\begin{itemize}
\item ABI version 5;
\item three bindings, A and B half and read-only, C float and written;
\item an empty constant pool;
\item 101 main instructions;
\item the 24 distinct (opcode, length) forms used;
\item the system registers read;
\item that exact launch geometry is required.
\end{itemize}

The Metal path turns this into an image (\cref{ch:image-result-validated}). The native path places the same bytes itself and writes the
launch words that correspond to these facts (\cref{ch:submission-below-metal}).

\section{General selection and allocation rules}\label{ch:selection-allocation-general}

As in the running example (\cref{ch:running-example-through}), \code{cc.select} uses a form only if its fields are mapped and its behaviour
was measured (the form registry, \code{isa/g17-forms.toml} and \code{isa/g17-scalar-isa.toml}). Anything else it refuses by name,
for example "no measured metadata layout for 5 bindings", or an unmeasured cooperative binding class
(\code{isa/g17-source-execution.json}). A refusal names the measurement that is missing.

Where the language leaves behaviour undefined, the compiler follows Apple's hardware. Its first lowering of
\code{(uint)x} for a negative float returned the magnitude, where Apple's code returns 0. It now lowers the cast to \op{9320}, reproducing Apple's bytes and output (MM~\mm{25.186}, MM~\mm{25.196}; \cref{sec:conversions}).

Allocation obeys \cref{ch:register-file}. \code{cc.allocate} assigns physical registers from R0 to R125 with tuple alignment, and
sets every lifetime modifier from liveness (\cref{sec:register-classes,sec:release-does}). It never releases a register
that is live around a back edge (\cref{sec:loops-loop-variable-rule}), it avoids encoding holes through the explicit
tables \code{\_encodable\_dests} and \code{\_encodable\_fma16\_sources} (\cref{sec:encoding-holes}), and it reserves every
register a tensor body names, through \code{\_prepare\_tensor\_composition} in \code{agxforge/g17/cc.py}.
\code{agxforge/g17/releasecheck.py} re-checks the emitted bytes for any read after a release.

\section{Tensor lowering}\label{ch:tensor-lowering}

\code{agxforge/g17/tlower.py} generates a matrix-product body from its shape and types (MM~\mm{23.4}, MM~\mm{25.165}):

\begin{enumerate}
\def\labelenumi{\arabic{enumi}.}
\item Read the lane index, and the simdgroup index when there are several. Compute row, column and element addresses
  in scalar code.
\item Build row and column masks for partial tiles. K beyond the matrix is zero-filled by the masked loads, never read
  out of range.
\item Allocate index scalars, fragments and eight-register accumulator groups.
\item Emit tensor loads with scoreboard slots and wait masks. The first K slice uses the from-zero MMA and later slices
  the accumulating one.
\item For a product that accumulates into C, compute from zero and add C afterwards, as Apple's library does.
\item Store the result, with masks on partial tiles.
\end{enumerate}

Larger shapes run as a runtime K loop that carries the accumulators across trips. It is bit-exact from 1 to 255
trips, and at M512 it is 2.25× faster than the unrolled body, because the unrolled form spills (MM~\mm{25.101}). Rule 30
of MM~\mm{17} is never to unroll the K loop.

A row grid, several simdgroups, a column grid and split-K spread a product across the GPU. Each combination is
admitted only where it was receipted on hardware (MM~\mm{25.134}, MM~\mm{25.165}).

Epilogues run inside the same dispatch. The output can be narrowed to half, requantized to int8 with
per-channel scales, or folded with a residual, without another dispatch (MM~\mm{25.165}, MM~\mm{25.183}).
